\section{Flow Grammar: Domain IR}
\label{sec:flow-grammar}

The Flow Grammar is the second grammar implemented in the diagram compiler and the \textbf{template grammar} for all future grammar additions. Where the Timeline Grammar (Section~\ref{sec:timeline-grammar}) proved the kernel's Scene IR contract with a time-axis-centric vocabulary, the Flow Grammar proves that the kernel is genuinely grammar-agnostic: it supports node-link topologies, pipeline sequences, and process flows---fundamentally different visual semantics---without kernel modifications.

% ============================================================================
\subsection{Scope \& Intent}
\label{sec:flow-scope}
% ============================================================================

\subsubsection{What a Flow Diagram IS}

A \textbf{flow diagram} in this grammar is a \emph{directed node-link diagram} representing:

\begin{itemize}
  \item \textbf{Pipelines/sequences}: ordered chains of processing steps (e.g., CI/CD, RAG retrieval, ETL).
  \item \textbf{Process flows}: multi-path directed graphs with branching and joining (e.g., decision trees, agent loops, state machines).
  \item \textbf{Architecture topologies}: layered or grouped node-link structures (e.g., microservice communication, data flow).
\end{itemize}

These correspond to the corpus taxonomy's \emph{node-link/topology} and \emph{pipeline/sequence} kinds (Section~\ref{sec:corpus-taxonomy}). The animated-arrow variant (flowing dashes along connectors) is a first-class pattern in this grammar.

\subsubsection{What a Flow Diagram is NOT}

\begin{description}
  \item[Not a comparison/matrix.] Tabular grid layouts with cells (Section~\ref{sec:corpus-comparison-matrix}) belong to the Comparison grammar. Flows have directed edges; comparisons do not.
  \item[Not a timeline.] The timeline grammar's entities have temporal semantics (\texttt{start}, \texttt{end}, \texttt{track}). Flow nodes have no time axis.
  \item[Not a general graph.] The Graph grammar (Section~\ref{sec:graph-grammar}) covers undirected networks, trees with Reingold-Tilford layout, and hub-and-spoke topologies. The Flow grammar is restricted to \emph{directed} structures with a declared layout intent---it does not attempt general force-directed placement.
\end{description}

\subsubsection{Modeling Boundary}

The Flow grammar covers diagrams where:
\begin{enumerate}
  \item Entities are discrete \textbf{nodes} (processing steps, decisions, data stores).
  \item Relationships are \textbf{directed edges} (data flow, control flow, sequence).
  \item Layout admits a natural \textbf{directional reading order} (left-to-right, top-to-bottom).
\end{enumerate}

Diagrams without directed edges, or without a dominant reading direction, are out of scope for this grammar.

% ============================================================================
\subsection{Flow Domain IR}
\label{sec:flow-domain-ir}
% ============================================================================

The Flow Domain IR is the small, grammar-specific representation that compiles \emph{down} to the shared Scene IR (Section~\ref{sec:scene-ir}). It is \textbf{not} a ``god-IR''---it contains only the constructs necessary to describe directed node-link diagrams with optional grouping.

\subsubsection{Document Structure}

A Flow IR document is a single YAML/JSON object with the following root structure:

\begin{lstlisting}[language=yaml,caption={Flow IR root document structure}]
version: "1.0"
metadata:
  title: string             # required
  subtitle: string          # optional; default null
  theme: string             # optional; default "default-flow"
flow:
  direction: left-to-right | top-to-bottom  # required
  nodes: [...]              # required; list<FlowNode>; may be empty
  edges: [...]              # optional; list<FlowEdge>; default []
  groups: [...]             # optional; list<FlowGroup>; default []
\end{lstlisting}

\subsubsection{Root Fields}

\begin{table}[h]
\centering
\small
\begin{tabular}{llccp{4.8cm}}
\toprule
\textbf{Field} & \textbf{Type} & \textbf{Req} & \textbf{Default} & \textbf{Description} \\
\midrule
\texttt{version} & \texttt{string} & Yes & --- & Spec version (semver, e.g., \texttt{"1.0"}) \\
\texttt{metadata} & \texttt{FlowMetadata} & Yes & --- & Document-level metadata \\
\texttt{flow} & \texttt{FlowDefinition} & Yes & --- & The flow graph definition \\
\bottomrule
\end{tabular}
\caption{Flow IR root fields}
\label{tab:flow-root-fields}
\end{table}

\subsubsection{FlowNode}
\label{sec:flow-node}

A node is a discrete entity in the flow: a processing step, decision point, data store, or any labeled box.

\begin{table}[h]
\centering
\small
\begin{tabular}{llccp{4.5cm}}
\toprule
\textbf{Field} & \textbf{Type} & \textbf{Req} & \textbf{Default} & \textbf{Description} \\
\midrule
\texttt{id} & \texttt{id} & Yes & --- & Document-unique identifier (kebab-case) \\
\texttt{label} & \texttt{string} & Yes & --- & Display text for the node \\
\texttt{shape} & \texttt{enum\{rect, rounded-rect, circle, diamond, stadium\}} & No & \texttt{rounded-rect} & Visual shape \\
\texttt{icon} & \texttt{string?} & No & \texttt{null} & Icon registry key \\
\texttt{status} & \texttt{enum\{default, active, success, warning, error, muted\}} & No & \texttt{default} & Semantic status (mapped to color by theme) \\
\texttt{description} & \texttt{string?} & No & \texttt{null} & Tooltip/secondary text \\
\texttt{group} & \texttt{ref<FlowGroup>?} & No & \texttt{null} & Group membership (alternative to declaring in group.nodes) \\
\bottomrule
\end{tabular}
\caption{FlowNode fields}
\label{tab:flow-node-fields}
\end{table}

\paragraph{Identity.}
Node \texttt{id} is the primary key within a document. It must be unique across all nodes. References in edges and groups use this \texttt{id}. Two nodes with the same \texttt{id} are a validation error.

\paragraph{Ordering.}
Nodes are ordered by their position in the \texttt{nodes} list. This order is the \textbf{canonical order} used by the layout engine for tie-breaking (determinism) and by serialization for stable output.

\subsubsection{FlowEdge}
\label{sec:flow-edge}

An edge represents a directed relationship between two nodes.

\begin{table}[h]
\centering
\small
\begin{tabular}{llccp{4.2cm}}
\toprule
\textbf{Field} & \textbf{Type} & \textbf{Req} & \textbf{Default} & \textbf{Description} \\
\midrule
\texttt{source} & \texttt{ref<FlowNode>} & Yes & --- & Source node id \\
\texttt{target} & \texttt{ref<FlowNode>} & Yes & --- & Target node id \\
\texttt{source\_port} & \texttt{enum\{auto, top, right, bottom, left\}} & No & \texttt{auto} & Anchor point on source node \\
\texttt{target\_port} & \texttt{enum\{auto, top, right, bottom, left\}} & No & \texttt{auto} & Anchor point on target node \\
\texttt{label} & \texttt{string?} & No & \texttt{null} & Edge annotation text \\
\texttt{style} & \texttt{enum\{solid, dashed, dotted\}} & No & \texttt{solid} & Stroke style \\
\texttt{animated} & \texttt{bool} & No & \texttt{false} & Enable flowing-dash animation \\
\texttt{directedness} & \texttt{enum\{directed, bidirectional, undirected\}} & No & \texttt{directed} & Arrow rendering \\
\bottomrule
\end{tabular}
\caption{FlowEdge fields}
\label{tab:flow-edge-fields}
\end{table}

\paragraph{Identity.}
Edges are identified by position in the \texttt{edges} list (index). Unlike nodes, edges do not carry an explicit \texttt{id}---their identity is their list position. This list position is the canonical order for deterministic processing.

\paragraph{Directedness.}
\begin{description}
  \item[\texttt{directed}] Arrow from source to target (default).
  \item[\texttt{bidirectional}] Arrows on both ends.
  \item[\texttt{undirected}] No arrowheads (connector line only).
\end{description}

\paragraph{Port resolution.}
When \texttt{source\_port} or \texttt{target\_port} is \texttt{auto}, the layout engine assigns the port based on the relative positions of source and target nodes. For left-to-right flows, \texttt{auto} resolves to \texttt{right} on source and \texttt{left} on target for forward edges.

\subsubsection{FlowGroup}
\label{sec:flow-group}

Groups are optional visual containers (lanes, clusters, swim-lanes) that partition nodes.

\begin{table}[h]
\centering
\small
\begin{tabular}{llccp{4.5cm}}
\toprule
\textbf{Field} & \textbf{Type} & \textbf{Req} & \textbf{Default} & \textbf{Description} \\
\midrule
\texttt{id} & \texttt{id} & Yes & --- & Document-unique identifier \\
\texttt{label} & \texttt{string} & Yes & --- & Group heading text \\
\texttt{nodes} & \texttt{list<ref<FlowNode>>} & No & \texttt{[]} & Node ids in this group (augments per-node \texttt{group} field) \\
\texttt{style} & \texttt{enum\{lane, cluster, outline\}} & No & \texttt{lane} & Visual rendering style \\
\bottomrule
\end{tabular}
\caption{FlowGroup fields}
\label{tab:flow-group-fields}
\end{table}

\paragraph{Group membership resolution.}
A node belongs to a group if \emph{either} it appears in the group's \texttt{nodes} list \emph{or} it has a \texttt{group} field referencing the group's \texttt{id}. Both mechanisms are equivalent; using both for the same node--group pair is valid (idempotent). A node may belong to at most one group; dual membership is a validation error.

\paragraph{Ungrouped nodes.}
Nodes not assigned to any group are rendered in the main flow area, outside any group container. They participate in layout normally.

\subsubsection{Layout Intent}
\label{sec:flow-layout-intent}

The \texttt{flow.direction} field declares the dominant layout direction:

\begin{description}
  \item[\texttt{left-to-right}] Nodes are ranked left-to-right; edges flow rightward. This is the default for pipeline/sequence diagrams.
  \item[\texttt{top-to-bottom}] Nodes are ranked top-to-bottom; edges flow downward. This suits vertical process flows.
\end{description}

The declared direction is a \textbf{hard constraint} on the layout engine: it determines layer-assignment orientation in the Sugiyama algorithm and reading order in the linear-sequence layout.

% ============================================================================
\subsection{Layout Contract}
\label{sec:flow-layout-contract}
% ============================================================================

The Flow grammar's layout engine converts the Domain IR into Scene IR coordinates. The choice of algorithm depends on graph topology, but all algorithms must satisfy the \textbf{determinism contract}: identical IR in $\rightarrow$ byte-identical Scene out.

\subsubsection{Layout Family Selection}

\begin{center}
\small
\begin{tabular}{lp{4.5cm}p{5.5cm}}
\toprule
\textbf{Topology} & \textbf{Algorithm} & \textbf{When Applied} \\
\midrule
Linear chain & Linear sequence layout & All nodes have in-degree $\le 1$ and out-degree $\le 1$ (a simple path or set of disjoint paths). \\
DAG & Sugiyama layered~\cite{sugiyama1981} & Directed acyclic graph with branching/joining. \\
Cyclic & Sugiyama with cycle removal & Cycles detected; feedback arcs reversed before layering. \\
\bottomrule
\end{tabular}
\end{center}

\paragraph{Topology detection is automatic.}
The layout engine inspects the graph structure and selects the appropriate algorithm. The author does not choose the algorithm---only the \texttt{direction}. This is a key difference from the Graph grammar (Section~\ref{sec:graph-grammar}), which requires explicit layout selection.

\subsubsection{Sugiyama Layered Layout (DAGs and Cyclic Flows)}

The primary layout algorithm for non-trivial flows follows the four-phase Sugiyama framework~\cite{sugiyama1981}:

\begin{enumerate}
  \item \textbf{Cycle removal.} Detect cycles via DFS. Reverse a minimal feedback arc set to produce a DAG. Reversed edges are rendered with a visual back-edge indicator (curved path routed behind the main flow).
  
  \item \textbf{Layer assignment.} Assign each node to a layer (rank) using network simplex~\cite{gansner1993dot}, minimizing total edge span. The \texttt{direction} field determines whether layers are columns (left-to-right) or rows (top-to-bottom).
  
  \item \textbf{Crossing minimization.} Barycenter heuristic iterated over adjacent layer pairs (24 sweeps, per standard practice~\cite{sugiyama1981}). Ties broken by canonical node list order (determinism).
  
  \item \textbf{Coordinate assignment.} Brandes \& K\"opf~\cite{brandesKopf2001} linear-time algorithm: four candidate alignments, median selected. Guarantees at most two bends per long edge.
\end{enumerate}

\paragraph{Edge routing.}
Edges are routed orthogonally where the Brandes--K\"opf coordinate assignment produces aligned segments. For non-orthogonal cases (long spans, back-edges), edges use smooth cubic B\'ezier curves. The choice is deterministic: given fixed node positions, edge routing is a pure function.

Orthogonal routing follows the principles of Tamassia's topology--shape--metrics framework~\cite{tamassia1987}: edges consist of horizontal and vertical segments only, with bends minimized by the coordinate-assignment phase. This produces clean, architecture-diagram-style connectors.

\subsubsection{Linear Sequence Layout (Pipelines)}

When the flow graph is a simple chain (every node has at most one predecessor and one successor), the layout engine uses a trivial linear placement:

\begin{itemize}
  \item Nodes placed at equal intervals along the primary axis (determined by \texttt{direction}).
  \item Edges are straight horizontal (LTR) or vertical (TTB) connectors.
  \item Groups (if present) span the extent of their member nodes.
\end{itemize}

This produces clean, minimal pipeline diagrams without the overhead of Sugiyama layer assignment.

\subsubsection{Determinism Mandate}

\textbf{All flow layouts are fully deterministic.} Specifically:

\begin{enumerate}
  \item No randomized initialization. No force-directed simulation. No RNG.
  \item Tie-breaking uses the canonical node order (list position in the IR).
  \item Crossing minimization uses a fixed sweep count (24), not convergence-based termination.
  \item The same IR document, processed by the same engine version, produces byte-identical Scene IR on every execution.
\end{enumerate}

If a force-like aesthetic is ever required for a flow diagram (not anticipated for v1), the implementation \textbf{must} use stress majorization~\cite{gansnerStressMaj2004} with a deterministic initial layout (nodes on a circle in canonical order) and a fixed iteration count. Raw Fruchterman-Reingold~\cite{fruchterman1991} with random initialization is prohibited.

\paragraph{Stress majorization rationale.}
Stress majorization (SMACOF) is a pure function of the graph-distance matrix given a fixed initial configuration. It converges monotonically (no oscillation), producing identical results across runs~\cite{gansnerStressMaj2004}. This is the only acceptable ``force-like'' family for a deterministic compiler.

% ============================================================================
\subsection{Lowering to the Scene IR}
\label{sec:flow-lowering}
% ============================================================================

The flow layout engine emits Scene IR primitives (Section~\ref{sec:scene-ir}). This section specifies the mapping from flow constructs to Scene primitives. \textbf{No new Scene IR primitives are required}---the existing kernel is sufficient.

\subsubsection{Node Lowering}

\begin{center}
\small
\begin{tabular}{lp{9cm}}
\toprule
\textbf{Flow Shape} & \textbf{Scene IR Primitive(s)} \\
\midrule
\texttt{rect} & \texttt{Rect\{corner\_radius: 0\}} \\
\texttt{rounded-rect} & \texttt{Rect\{corner\_radius: theme.flow.nodeRadius\}} \\
\texttt{circle} & \texttt{Circle} \\
\texttt{diamond} & \texttt{Path} (four-point rhombus via \texttt{M/L/L/L/Z} commands) \\
\texttt{stadium} & \texttt{Rect\{corner\_radius: height/2\}} (pill shape) \\
\bottomrule
\end{tabular}
\end{center}

Each node emits a \texttt{Group} containing:
\begin{enumerate}
  \item The shape primitive (background/border).
  \item A \texttt{Text} primitive for the label (centered within the shape).
  \item Optionally, an \texttt{Image} primitive for the icon (positioned above or left of the label, per theme).
\end{enumerate}

\paragraph{Status $\rightarrow$ color.}
The node's \texttt{status} field is resolved by the theme to \texttt{fill} and \texttt{stroke} colors on the shape primitive. The Domain IR does \emph{not} carry color values---color is always theme-resolved.

\subsubsection{Edge Lowering}

Each edge emits:

\begin{enumerate}
  \item A \texttt{Path} primitive for the connector line (routed by the layout engine). Stroke style (\texttt{solid}, \texttt{dashed}, \texttt{dotted}) maps directly to the Scene Path's \texttt{stroke\_style} field.
  
  \item An arrowhead: a small filled \texttt{Path} (triangle) at the target end for \texttt{directed} edges; at both ends for \texttt{bidirectional}; omitted for \texttt{undirected}.
  
  \item Optionally, a \texttt{Text} primitive for the edge label, positioned at the midpoint of the path (offset perpendicular to the edge direction to avoid overlap).
\end{enumerate}

\paragraph{Animated edges.}
When \texttt{animated: true}, the edge's \texttt{Path} primitive receives a \texttt{FlowingDashes} animation hint (Section~\ref{sec:scene-animation}):

\begin{Verbatim}[frame=single,fontsize=\small]
FlowingDashes {
  target_id: <edge_path_id>,
  dash_length: theme.flow.animDashLength,     // default 8px
  gap_length:  theme.flow.animGapLength,      // default 4px
  speed_px_per_sec: theme.flow.animSpeed,     // default 40
  direction: 'forward'
}
\end{Verbatim}

SVG/HTML backends render this as \texttt{stroke-dashoffset} animation. Raster backends (PNG/Skia) render the resting frame: a static dashed stroke.

\subsubsection{Group Lowering}

Groups emit:
\begin{enumerate}
  \item A background \texttt{Rect} spanning the bounding box of all member nodes (with padding from theme).
  \item A \texttt{Text} primitive for the group label (positioned at the top of the rect).
  \item All member nodes rendered as children of a \texttt{Group} primitive (for logical containment, not visual nesting---nodes are positioned absolutely).
\end{enumerate}

\paragraph{No new primitives needed.}
The flow grammar maps entirely to existing Scene IR primitives: \texttt{Rect}, \texttt{Circle}, \texttt{Path}, \texttt{Text}, \texttt{Image}, and \texttt{Group}. No kernel changes are required.

\textbf{Open question for Barbara/Mark:} If future flow diagrams require \emph{nested} groups (groups within groups), should the Scene IR \texttt{Group} primitive support recursive clip regions, or should nesting be flattened at lowering time? Current spec assumes single-level groups only. Nested groups are deferred to v2.

% ============================================================================
\subsection{Edge Cases \& Total Semantics}
\label{sec:flow-edge-cases}
% ============================================================================

Every possible input must have a defined meaning or produce a clear validation error. This section pins down the edge cases.

\subsubsection{Cycles}

\textbf{Cycles are valid.} A flow may contain cycles (e.g., retry loops, feedback arcs in agent pipelines). The layout engine handles cycles by:

\begin{enumerate}
  \item Detecting strongly-connected components via Tarjan's algorithm.
  \item Selecting a minimal feedback arc set (edges to reverse for layering). Selection is deterministic: among equally-scored candidates, the edge appearing \emph{later} in the canonical edge list is chosen for reversal.
  \item Reversed edges are rendered as \emph{back-edges}: routed with a longer path (curved or stepped) that visually indicates ``going back'' in the flow direction.
\end{enumerate}

The resulting diagram remains readable: forward edges flow in the declared direction; back-edges are visually distinct.

\subsubsection{Self-Loops}

\textbf{Self-loops are valid.} An edge where \texttt{source == target} is rendered as a loopback connector: a curved path that exits the node from one port and re-enters at an adjacent port (e.g., exits \texttt{right}, loops around, enters \texttt{top}). The specific port pair is determined by the flow direction and node position.

\subsubsection{Multi-Edges}

\textbf{Multiple edges between the same pair of nodes are valid.} Each edge is rendered as a separate path, offset perpendicular to the primary path direction to avoid visual overlap. Offset distance is \texttt{theme.flow.multiEdgeGap} (default: 6px). Edges are offset in canonical order (list position).

\subsubsection{Disconnected / Orphan Nodes}

\textbf{Nodes with no edges are valid.} Orphan nodes (zero in-degree and zero out-degree) are placed in a designated ``orphan region''---below the main flow (for LTR) or to the right (for TTB)---sorted by canonical order. They participate in group membership normally.

\textbf{Disconnected components} (multiple disjoint subgraphs) are each laid out independently using the appropriate algorithm, then arranged sequentially along the secondary axis (vertical for LTR, horizontal for TTB).

\subsubsection{Missing Target/Source ID}

\textbf{Validation error.} If an edge references a node \texttt{id} that does not exist in the \texttt{nodes} list, the document fails validation with a diagnostic: \texttt{"Edge at index N references unknown node id 'X' in field 'source|target'"}. The layout engine does not process invalid documents.

\subsubsection{Empty Flow}

\textbf{Valid but degenerate.} A flow with zero nodes and zero edges produces a Scene with only the canvas and metadata (no drawing primitives). This is analogous to the timeline grammar's empty-activities case. A lint warning is emitted: \texttt{"Flow contains no nodes; output will be empty."}.

\subsubsection{Duplicate Node IDs}

\textbf{Validation error.} Two nodes with the same \texttt{id} fail validation: \texttt{"Duplicate node id 'X' at indices M and N."}.

\subsubsection{Edge Referencing Same Source and Target (Self-Loop)}

Covered above (\S self-loops). This is valid, not an error.

% ============================================================================
\subsection{Worked Example: RAG Pipeline}
\label{sec:flow-worked-example}
% ============================================================================

This example is drawn from the corpus (rag-vectorless-explainer, Section~\ref{sec:corpus-taxonomy}): a simplified Retrieval-Augmented Generation pipeline.

\subsubsection{Flow Domain IR}

\begin{lstlisting}[language=yaml,caption={RAG pipeline --- Flow Domain IR}]
version: "1.0"
metadata:
  title: "RAG Retrieval Pipeline"
  theme: dark-flow
flow:
  direction: left-to-right
  nodes:
    - id: query
      label: "User Query"
      shape: stadium
      icon: message-circle
      status: active
    - id: embed
      label: "Embed Query"
      shape: rounded-rect
      icon: cpu
    - id: retrieve
      label: "Vector Search"
      shape: rounded-rect
      icon: database
    - id: rerank
      label: "Re-rank"
      shape: rounded-rect
      icon: filter
    - id: generate
      label: "Generate Answer"
      shape: rounded-rect
      icon: sparkles
      status: success
  edges:
    - source: query
      target: embed
      style: solid
    - source: embed
      target: retrieve
      style: solid
      animated: true
    - source: retrieve
      target: rerank
      style: dashed
      label: "top-k"
    - source: rerank
      target: generate
      style: solid
      animated: true
    - source: generate
      target: query
      style: dotted
      label: "follow-up"
      directedness: directed
  groups:
    - id: retrieval
      label: "Retrieval Stage"
      nodes: [embed, retrieve, rerank]
      style: lane
\end{lstlisting}

\subsubsection{Lowering to Scene IR (Prose Description)}

The layout engine processes this document as follows:

\begin{enumerate}
  \item \textbf{Topology detection.} The edge \texttt{generate $\rightarrow$ query} creates a cycle. The engine identifies this as a feedback arc and reverses it for layering purposes.

  \item \textbf{Layer assignment.} Network simplex assigns five layers (left-to-right): L0=\texttt{query}, L1=\texttt{embed}, L2=\texttt{retrieve}, L3=\texttt{rerank}, L4=\texttt{generate}. The reversed back-edge \texttt{generate $\rightarrow$ query} spans 4 layers.

  \item \textbf{Crossing minimization.} With a single node per layer, no crossings exist. The barycenter pass completes trivially.

  \item \textbf{Coordinate assignment.} Nodes are placed at equal horizontal intervals. The ``Retrieval Stage'' group's bounding box spans L1--L3.

  \item \textbf{Scene IR emission.} The resulting Scene contains:
  \begin{itemize}
    \item 5 \texttt{Group} primitives (one per node), each containing a \texttt{Rect} (stadium/rounded-rect) + \texttt{Text} (label) + \texttt{Image} (icon).
    \item 5 \texttt{Path} primitives (edges). The \texttt{embed$\rightarrow$retrieve} and \texttt{rerank$\rightarrow$generate} paths carry \texttt{FlowingDashes} animation hints. The \texttt{generate$\rightarrow$query} back-edge is routed as a curved path below the main flow.
    \item 5 small \texttt{Path} primitives (arrowheads on directed edges).
    \item 2 \texttt{Text} primitives for edge labels (``top-k'' and ``follow-up'').
    \item 1 \texttt{Rect} (group background for ``Retrieval Stage'') + 1 \texttt{Text} (group label).
    \item Canvas: dimensions computed from content bounding box + padding.
    \item Meta: \texttt{scene\_hash} computed from the serialized element list.
  \end{itemize}
\end{enumerate}

\paragraph{Determinism verification.}
Running this IR through the layout engine twice produces an identical \texttt{scene\_hash}. The cycle-handling, layer assignment, and coordinate assignment are all pure functions of the input and canonical ordering.

% ============================================================================
\subsection{Determinism \& Opt-In Discipline}
\label{sec:flow-determinism}
% ============================================================================

The flow grammar respects the project's sacred determinism contract (Section~\ref{sec:determinism}):

\begin{description}
  \item[Byte-identical output.] The same Flow IR + same engine version + same theme = identical Scene IR (and therefore identical SVG). No exceptions.
  
  \item[No-change defaults.] All optional fields (\texttt{shape}, \texttt{status}, \texttt{style}, \texttt{animated}, \texttt{directedness}, ports, groups) have explicit defaults. Omitting optional fields produces the same output as specifying the default value explicitly.
  
  \item[New tokens are opt-in.] Future additions to the Flow IR (e.g., new shape values, new edge styles) must have no-change defaults: existing documents that do not use new features must produce byte-identical output after an engine upgrade.
  
  \item[Canonical ordering.] List position in \texttt{nodes} and \texttt{edges} is the tie-breaking order for all layout decisions. Reordering nodes in the IR may change layout (this is intentional---order is semantic).
  
  \item[Animation does not break determinism.] The \texttt{animated} flag adds a declarative \texttt{FlowingDashes} hint to the Scene. The SVG markup is byte-identical regardless of whether the animation is ``playing'' --- animation is CSS/SMIL, not frame-by-frame.
\end{description}

% ============================================================================
\subsection{Open Questions (Deferred to Mark/Barbara)}
\label{sec:flow-open-questions}
% ============================================================================

The following design points are flagged for specialist review before the schema is finalized:

\begin{description}
  \item[Mark (IR schema detail):]
  \begin{itemize}
    \item Exact JSON Schema for the Flow Domain IR (enum values, string patterns, min/max constraints).
    \item Whether \texttt{FlowEdge} should carry an explicit \texttt{id} field for referencing in composition or annotation contexts.
    \item Port model: is the 5-value enum (\texttt{auto, top, right, bottom, left}) sufficient, or should named custom ports be supported (useful for nodes with multiple outgoing paths)?
    \item Validation rule set: exhaustive list of semantic validation rules beyond schema conformance.
  \end{itemize}
  
  \item[Barbara (rendering semantics):]
  \begin{itemize}
    \item Self-loop routing: exact curve parameters (radius, port pair selection heuristic).
    \item Back-edge rendering: curve style (cubic B\'ezier, stepped, or arc) and visual offset distance.
    \item Multi-edge offset geometry: perpendicular offset vs.\ curvature-based separation.
    \item Group visual style details: border radius, padding, header positioning rules.
    \item Edge label collision avoidance with nodes/other labels: post-layout adjustment rules.
  \end{itemize}
\end{description}

\paragraph{Kernel change flag.}
No kernel (Scene IR) changes are required for the flow grammar as specified. If nested groups or custom port models later prove necessary, those would require kernel-level changes needing Barbara and Mark sign-off before adoption.
